\documentclass[11pt,a4paper]{article}
\usepackage[T1]{fontenc}
\usepackage[utf8]{inputenc}
\usepackage{lmodern}
\usepackage[margin=27mm,headheight=15pt]{geometry}
\usepackage{amsmath,amssymb,amsthm,mathtools}
\usepackage{booktabs,longtable,array,enumitem,microtype}
\usepackage{xcolor,fancyhdr}
\usepackage[unicode,colorlinks=true,linkcolor=blue!45!black,citecolor=blue!45!black,urlcolor=blue!45!black]{hyperref}
\usepackage{bookmark}
\definecolor{ink}{RGB}{25,42,64}
\definecolor{accent}{RGB}{35,91,120}
\newcommand{\R}{\mathbb R}
\newcommand{\N}{\mathbb N}
\newcommand{\E}{\mathbb E}
\newcommand{\Pbb}{\mathbb P}
\newcommand{\norm}[1]{\left\lVert#1\right\rVert}
\newcommand{\abs}[1]{\left\lvert#1\right\rvert}
\newcommand{\dd}{\,\mathrm d}
\DeclareMathOperator{\dist}{dist}
\DeclareMathOperator{\supp}{supp}
\DeclareMathOperator{\spanop}{span}
\DeclareMathOperator{\diam}{diam}
\DeclareMathOperator{\Lip}{Lip}
\newtheorem{theorem}{Theorem}[section]
\newtheorem{lemma}[theorem]{Lemma}
\newtheorem{proposition}[theorem]{Proposition}
\newtheorem{corollary}[theorem]{Corollary}
\theoremstyle{definition}
\newtheorem{definition}[theorem]{Definition}
\theoremstyle{remark}
\newtheorem{remark}[theorem]{Remark}
\numberwithin{equation}{section}
\setlist{topsep=4pt,itemsep=2pt,parsep=0pt}
\setlength{\emergencystretch}{2em}
\pagestyle{fancy}
\fancyhf{}
\fancyhead[L]{\small\color{ink}AIM 568}
\fancyhead[R]{\small 4 November 2026}
\fancyfoot[C]{\thepage}
\renewcommand{\headrulewidth}{0.3pt}


\urlstyle{same}
\title{A linear information bound for nonsymmetric convex classes under scalar-error measurability}
\author{Matthew J. Colbrook\thanks{The solutions were found by AI. The submitted arguments and the repository audit are AI work; independent human review and formal verification are not claimed.}\\[4pt]
\small Department of Applied Mathematics and Theoretical Physics\\
\small University of Cambridge, Cambridge CB3 0WA, United Kingdom\\
\small\href{mailto:m.colbrook@damtp.cam.ac.uk}{m.colbrook@damtp.cam.ac.uk}}
\date{4 November 2026}
\hypersetup{pdftitle={A linear information bound for nonsymmetric convex classes under scalar-error measurability},pdfauthor={Matthew J. Colbrook},pdfsubject={AIM 568: AI-found solution or partial result}}
\begin{document}
\maketitle
\begin{abstract}
We extend the $27n$ comparison of Krieg, Novak and Ullrich to the scalar-error measurability convention in AIM 568. A finite-dimensional measurable-majorant lemma supplies the bridge without assuming that the original adaptive rules or reconstruction are measurable. The source theorem supplies the geometric lower bound; the contribution here is this bridge.
\end{abstract}
\noindent\textbf{Catalogue target.} \href{https://github.com/MColbrook/AIM/blob/c4036871b0bf15cbc83e79c780b7191bcbb2cfed/problems/568-a-linear-information-bound-for-arbitrary-convex-model-classes.md}{AIM 568 at repository revision \texttt{c4036871b0bf}}. The package preserves its full statement in \texttt{statement.md}; the review and target comparison are linked from the package README.
\medskip

\section{Problem 568: closing the measurability gap}
\label{sec:568}

Throughout, \(e_n^{\rm ran}(S,F)\) uses expected norm error and a strict
budget of at most \(n\) queries in every realization, as required by AIM 568.
The class \(F\subset X\) is bounded and convex, without a symmetry
assumption. Its measurability convention uses the norm
\(p_F\) on \(\operatorname{span}(F-F)\).

Krieg, Novak and Ullrich's October 2026 preprint
\cite{knu2026} proves the \(27n\) comparison with measurable adaptive
decision rules and reconstruction. The catalogue retains only
measurability of the scalar error, so the theorem cannot simply be
quoted in its stated model. The result below supplies the missing
bridge.

\begin{theorem}
\label{thm:568}
Under exactly the assumptions of Problem 568,
\[
 e_{2n}^{\rm det,non}(S,F)
 \le e_{2n-1}^{\rm det,non}(S,F)
 \le a_{2n-1}(S,F)
 \le 27n\,e_n^{\rm ran}(S,F),\qquad n\ge1,
 \tag{568.1}
\]
where \(a_m\) is the minimal worst-case error of continuous affine
algorithms whose linear part has rank at most \(m\).
\end{theorem}

The numerical comparison constant is inherited from
\cite{knu2026}. The contribution here is its extension to the larger
admissibility class specified in the question.

\subsection{The catalogue measurability convention}
Let $X,Y$ be real Banach spaces, $S:X\to Y$ bounded and linear, and $F\subset X$ nonempty, bounded and convex. On $\operatorname{span}(F-F)$ define
\[
p_F(v)=\inf\{t>0:v\in t(F-F)\}.
\]
The input Borel structure is that of the metric $d_F(f,g)=p_F(f-g)$ (the one-point topology in the trivial case). Each randomized realization uses at most $n$ adaptive continuous linear measurements. Only the scalar error $(f,\omega)\mapsto\|Sf-A_\omega(f)\|_Y$ must be jointly measurable. Its criterion is $\sup_{f\in F}\mathbb E\|Sf-A_\omega(f)\|_Y$, denoted $e_n^{\rm ran}(S,F)$ after minimization. The deterministic quantity $e_m^{\rm det,non}(S,F)$ minimizes the worst-case norm error over $m$ fixed continuous linear measurements and arbitrary reconstruction. We prove the requested comparison with $C=27$, in fact at budget $2n-1$.

\subsection{A measurable-majorant principle}

\begin{lemma}
\label{lem:majorant}
Let \(K\subset\R^d\) be compact, \(T:K\to\R^q\) continuous,
\(\nu\) a Borel probability measure on \(K\), and
\(g:K\to[0,\infty]\) Borel. Suppose an algorithm \(A\), with entirely
arbitrary adaptive decisions and decoder, uses at most \(k\) affine
scalar measurements and satisfies
\[
                       \norm{T(x)-A(x)}_2\le g(x)
 \tag{B1}
\]
outside a Borel \(\nu\)-null set. Then, for every \(\varepsilon>0\),
there is an algorithm \(B\) with at most \(k\) affine scalar
measurements, Borel history-to-query maps and a Borel decoder, such
that
\[
       \int_K\norm{T(x)-B(x)}_2\,d\nu(x)
                    \le \int_K g\,d\nu+\varepsilon .
 \tag{B2}
\]
The algorithm's own error in (B1) need not be measurable.
\end{lemma}

\begin{proof}
Choose a closed Euclidean ball \(C\) of radius \(R\) containing \(T(K)\).
Projecting a decoder onto \(C\) cannot increase its distance to any
\(T(x)\). We may therefore restrict all candidate decoders to \(C\);
their losses are bounded by \(2R\).

For each \(j\ge0\), let \(\mathcal C_j\) be the countable family of
\(j\)-query policies whose affine query coefficients are
rational-coefficient polynomials of the previous answers and whose
decoder is the projection onto \(C\) of a rational-coefficient
polynomial map. Each complete input-output map in this family is
continuous on \(K\).

We first establish, for every Borel probability \(\lambda\) on \(K\),
\[
 \inf_{B\in\mathcal C_j}\int_K\norm{T-B}_2\,d\lambda
   =\inf_{\substack{B:\ j\text{ queries}\\
                    \text{Borel decisions and decoder}}}
                  \int_K\norm{T-B}_2\,d\lambda .
 \tag{B3}
\]
Start with a Borel policy and consider each of its \emph{original}
finite-history probability laws. By Lusin's theorem, its query
coefficient map is continuous on a compact set of that history space
having probability at least \(1-\delta\). Extend the coefficients
continuously by coordinatewise Tietze extension. Do the same for the
decoder, then project the extension onto \(C\). On the event that
every original history belongs to its designated compact set, the
modified policy follows precisely the same history and gives the
same output, by induction through the stages. The complement has
probability at most \((j+1)\delta\). Hence the change in mean loss is
at most \(2R(j+1)\delta\).

The modified policy is continuous in its histories. Its reachable
histories on compact \(K\) are compact. Approximate all coefficient
maps and the decoder on compact boxes containing neighborhoods of
those histories by rational-coefficient polynomials. Successive
uniform continuity through the finite observation recursion shows
that sufficiently close approximations give uniformly close final
outputs on \(K\). Projection onto \(C\) is nonexpansive. This proves
(B3); the reverse inequality is immediate.

Define
\[
              V_j(\lambda)=
              \inf_{B\in\mathcal C_j}
                    \int_K\norm{T(x)-B(x)}_2\,d\lambda(x).
 \tag{B4}
\]
For each fixed \(B\in\mathcal C_j\), its loss is continuous and
bounded on \(K\), so its integral is continuous for the weak topology
on \(\mathcal P(K)\). Thus \(V_j\) is an infimum of countably many
continuous functions, and is upper semicontinuous and Borel.
By (B3), it is exactly the optimal Borel-policy average error.

We now prove \(\int g\,d\nu\ge V_k(\nu)\) by induction on \(k\).
For \(k=0\), \(A\) is constant, and the claim follows by projecting
its value to \(C\). For \(k\ge1\), fix the first affine query and
subtract its known constant, leaving \(x\mapsto a\cdot x\).
There is a Borel regular conditional-law kernel
\[
                  y\longmapsto\nu_y\in\mathcal P(K)
 \tag{B5}
\]
for \(\nu\) given \(a\cdot x=y\), supported on that fiber for almost
every \(y\). This standard conditional law can also be constructed
directly: condition on the dyadic interval of length \(2^{-m}\)
containing \(y\), take normalized restrictions, and then take their
weak limit where it exists. Martingale convergence for a countable
dense family in \(C(K)\), together with compactness of
\(\mathcal P(K)\), gives that limit almost everywhere; use a fixed
point mass elsewhere. This yields a Borel kernel.

For almost every \(y\), the Borel null witness in (B1) has
\(\nu_y\)-measure zero. Freeze the first history. The remaining
\(k-1\) stages are some well-defined adaptive algorithm on the fiber,
whether or not their choice depends measurably on \(y\). Induction
applied separately to this conditional law gives
\[
                        \int_K g\,d\nu_y\ge V_{k-1}(\nu_y).
 \tag{B6}
\]
Enumerate \(\mathcal C_{k-1}=(B_i)_{i\ge1}\). For each \(y\), choose
the least \(i=i(y)\) satisfying
\[
       \int_K\norm{T-B_i}_2\,d\nu_y
                          \le V_{k-1}(\nu_y)+\varepsilon .
 \tag{B7}
\]
All quantities in (B7) are Borel in \(y\), so this is a Borel
integer-valued choice. The policy which first measures \(a\cdot x\)
and then follows \(B_{i(y)}\) is a Borel \(k\)-query policy. Ordinary
disintegration of its bounded Borel loss gives
\[
\begin{split}
 V_k(\nu)
 &\le \int_\R V_{k-1}(\nu_y)\,d(a_*\nu)(y)+\varepsilon\\
 &\le \int_K g\,d\nu+\varepsilon .
\end{split}
 \tag{B8}
\]
Let \(\varepsilon\downarrow0\), and then use (B3) once more to obtain
(B2). The proof has never integrated the original nonmeasurable
query choices or decoder.
\end{proof}

\subsection{Apply the bridge to the fixed simplex lower bound}

\begin{proof}[Proof of Theorem~\ref{thm:568}]
We specify exactly which geometric facts from \cite{knu2026} are
used. For \(j\ge1\), define
\[
 D_j=\sup_{\substack{f_0,\ldots,f_j\in F\\
                    Q:Y\to\ell_2^j,\ \norm Q\le1}}
       \abs{\det[QS(f_1-f_0),\ldots,QS(f_j-f_0)]}.
 \tag{568.2}
\]
If \(D_{2n}=0\), the affine span of \(S(F)\) has dimension at most
\(2n-1\). A bounded projection onto that finite-dimensional
direction space then gives an exact continuous affine algorithm,
and (568.1) is immediate.

Suppose \(D_{2n}>0\), and write
\[
                         Q_n=(D_{2n}/D_n)^{1/n}.
 \tag{568.3}
\]
Lemma 2.2 of \cite{knu2026} gives
\[
                      a_{2n-1}(S,F)\le2\sqrt n\,Q_n.
 \tag{568.4}
\]
Its Lemma 2.3, Proposition 3.2, and the constant calculation in
Sections 4--5 give the following stronger ingredient underlying
the randomized theorem. For every \(0<\eta<1\), one can fix a
\(2n\)-dimensional simplex \(\Delta\subset F\), a contraction
\(Q:Y\to\R^{2n}\), and an affine parametrization
\[
             f=\bar f+Vz,\quad z\in K,\qquad T=QSV,
 \tag{568.5}
\]
where \(K\) is a compact centered simplex with its uniform
probability law \(\nu\), such that every Borel finite-dimensional
algorithm \(\widetilde B\) using at most \(n\) affine measurements
satisfies
\[
 \int_K\norm{Tz-\widetilde B(z)}_2\,d\nu(z)
      \ge L_\eta,\qquad
 L_\eta=\frac{2(1-\eta)^{1/n}}{27\sqrt n}\,Q_n .
 \tag{568.6}
\]
The simplex, contraction and probability law are chosen independently
of the algorithm and its random seed. Equation (568.6) is the
Euclidean uniform-simplex lower bound from that proof, not an
application of its global randomized theorem to a different model.

Now take an algorithm admissible in the catalogue's broader model.
For each seed \(\omega\), put
\[
\begin{split}
 g_\omega(z)
   &=\norm{S(\bar f+Vz)-A_\omega(\bar f+Vz)}_Y,\\
 \widetilde A_\omega(z)
   &=Q A_\omega(\bar f+Vz)-QS\bar f .
\end{split}
 \tag{568.7}
\]
The function \(g_\omega\) is Borel. To check the precise topology,
the restriction of \(p_F\) to the finite-dimensional span of the
simplex differences is a norm, hence is equivalent to the Euclidean
norm. Thus \(z\mapsto\bar f+Vz\) is continuous for the stipulated
input topology. Joint measurability in \((z,\omega)\) follows in
the same way.

The algorithm \(\widetilde A_\omega\) uses at most \(n\) affine
scalar measurements on \(K\), and contraction gives the pointwise
majorization
\[
               \norm{Tz-\widetilde A_\omega(z)}_2
                              \le g_\omega(z).
 \tag{568.8}
\]
The left side need not be measurable. This is why a
\emph{measurable-majorant} lemma, rather than an unqualified
measurable-output replacement, was proved.

Apply Lemma~\ref{lem:majorant} to (568.8). Since every Borel simulator
is bounded below by (568.6), it follows that
\[
                         \int_K g_\omega\,d\nu\ge L_\eta
                          \quad\text{for every }\omega.
 \tag{568.9}
\]
Tonelli, using only the scalar error's given joint measurability,
therefore yields
\[
 \sup_{f\in F}\E_\omega\norm{Sf-A_\omega(f)}_Y
       \ge \int_K\E_\omega g_\omega(z)\,d\nu(z)
       \ge L_\eta .
 \tag{568.10}
\]
Take the infimum over all catalogue-admissible algorithms, let
\(\eta\downarrow0\), and combine with (568.4):
\[
       a_{2n-1}(S,F)\le27n\,e_n^{\rm ran}(S,F).
 \tag{568.11}
\]
Finally, an affine algorithm of rank at most \(2n-1\) is realized
by that many fixed continuous linear measurements. Monotonicity
in the deterministic budget proves all inequalities in (568.1).
\end{proof}


\begin{thebibliography}{99}
\bibitem{knu2026}D. Krieg, E. Novak and M. Ullrich, \emph{On the power of adaption and randomization 2: Non-symmetric convex sets}, 24 October 2026, arXiv:2610.33733v1, especially Theorem 1.1, Proposition 3.2 and its geometric proof. \url{https://arxiv.org/abs/2610.33733}.
\end{thebibliography}
\end{document}
